\documentclass[12pt]{report}
\usepackage[usenames,dvipsnames]{color}
\usepackage{listings}
\usepackage{graphicx}
\usepackage{fancyhdr}
\usepackage{framed}
\usepackage[T1]{fontenc}
\usepackage[toc,page]{appendix}
\usepackage[utf8]{inputenc}
\usepackage[brazil]{babel}
\usepackage{fancyvrb}
\usepackage[hmargin=2cm,vmargin=2cm]{geometry}
\usepackage{lastpage}
\usepackage{pdfpages}
\usepackage{makeidx}
\usepackage{hyperref}
\pagestyle{fancy}
\usepackage{enumitem}
% cabecalho e rodapé
\setlength{\headheight}{120pt}
\setlength{\textheight}{550pt}
\renewcommand{\headrulewidth}{0pt}
\lhead{\includegraphics[scale=0.03]{brasao.png}}
\chead{\includegraphics[scale=0.5]{logo-brasil-sem-pobreza2.png}}
\rhead{\includegraphics[scale=0.5]{logo-pnud.png}}
\cfoot{\textbf{\ProjectCode\ - Inovando a democracia participativa}}
\rfoot{\thepage}

\hyphenation{par-ti-ci-pa-ção}
\bibliographystyle{ieeetr}

% definições sobre o autor e o produto
\newcommand{\MyName}{Renato Fabbri}
\newcommand{\MySurnameForename}{Fabbri, Renato}
\newcommand{\SupervisorName}{Gabriella Vieira Oliveira Gonçalves}
\newcommand{\MyEmail}{renato.fabbri@gmail.com}
\newcommand{\ContractNumber}{2013/000566}
\newcommand{\ContractYear}{2014}
\newcommand{\ProjectCode}{Projeto BRA/12/018}
\newcommand{\NomeSecretaria}{Secretaria-Geral da Presidência da República}
%Q\newcommand{\SiglaSecretaria}{SG/PR}
\newcommand{\SiglaSecretaria}{Secretaria: SNAS }
\newcommand{\ProductNumber}{Extra}
\newcommand{\ProductTitle}{Anotações sobre a leitura de produtos dos outros consultores do Projeto BRA/12/018}
\newcommand{\ProductSubtitle}{confluências, aplicações, dúvidas, sugestões, correções}
\newcommand{\ProductDescription}{"Ferramentas assistidas de categorização de conteúdo: Com Processamento de Linguagem Natural e de Redes Complexas, adaptadas para o ambiente do portal de participação."
}

\newcommand{\ProductValue}{R\$ 0,00 (zero reais e zero centavos)}
\newcommand{\ObjetoContratacao}{
Aporte de conhecimentos e tecnologias para especificação de vocabulário e ferramentas assistidas que utilizam processamento de linguagem natural e análise de redes complexas para o conteúdo do portal da participação social.
}
\newcommand{\DataEntrega}{Agosto de 2014}
\newcommand{\PalavrasChave}{reconhecimento de padrões, redes complexas, processamento de linguagem natural, participação social}
\newcommand{\pp}[1]{

\textbf{Parágrafo #1:}

}
% lista de abreviações
\makeindex

\begin{document}

\newgeometry{hmargin=3cm,vmargin=1.5cm}
\begin{center}
\thispagestyle{empty}
{\color{MidnightBlue}

\includegraphics[scale=0.9]{logo-pnud.png}

\vspace{4cm}

{\bf \large \ProjectCode\ - Desenvolvimento de Metodologias
de Articulação e Gestão de Políticas Públicas para Promoção da Democracia
Participativa}

\vspace{1.5cm}

{\bf \large Produto \ProductNumber\ -\ \ProductTitle}

\vspace{1.5cm}

\ProductSubtitle

\vspace{4cm}

\MyName

\vspace{2cm}

}

\includegraphics[scale=0.04]{brasao.png} \\
{\bf \small \NomeSecretaria}

\end{center}
\restoregeometry
\newpage

\newgeometry{hmargin=3cm,vmargin=1.5cm}
\addtolength{\topmargin}{2.5cm}
\thispagestyle{empty}
{\color{MidnightBlue}

{\bf \LARGE Produto \ProductNumber\ -\ \ProductTitle}

\hrulefill

\vspace{1cm}

\begin{center}

{\bf \large Contrato n. \ContractNumber}

\vspace{1.5cm}

{\bf \large Objeto da contratação: \ObjetoContratacao}

\end{center}

\vspace{3.2cm}

Valor do produto: \ProductValue

\vspace{1.2cm}

Data de entrega: \DataEntrega

\vspace{1.2cm}

Nome do consultor: \MyName

\vspace{1.2cm}

Nome da supervisora: \SupervisorName

}

\vspace{2cm}

\begin{center}
\includegraphics[scale=0.04]{brasao.png} \\
{\bf \small \NomeSecretaria}
\end{center}

\restoregeometry
\newpage

\newgeometry{hmargin=3cm,vmargin=1.5cm}
\addtolength{\topmargin}{5cm}
\thispagestyle{empty}

\begin{framed}

{\raggedright \MySurnameForename} \\

\ProductTitle: \ProductSubtitle\ / \ContractYear. \\

Total de folhas: \pageref{LastPage} \\

\vspace{1cm}

Supervisora: \SupervisorName \\

\SiglaSecretaria \\

\NomeSecretaria \\

Palavras-chave: \PalavrasChave. \\

\end{framed}

\vspace{3cm}

{\raggedright \includegraphics{licenca-cc-by-nc.png} \ Esta obra é licenciada sob
uma licença Creative Commons - Atribuição-NãoComercial. 4.0 Internacional.}

\restoregeometry
\newpage

\tableofcontents
\newpage


\begin{abstract}
Este documento registra a reflexão sobre produtos dos outros consultores do mesmo projeto (BRA/12/018).
Cada consultor entrega alguns ``produtos'' para os órgãos interessados. No caso, são documentos escritos,
que relatam atividades, pesquisas, propostas, enfim, o que for pertinente para o trabalho.
Na leitura destes documentos, são registradas anotações pertinentes à consultoria de contrato 2013/000566, sob responsabilidade
do autor deste produto extra. Cada parte do documento corresponde a um consultor, cada capítulo corresponde
a um documento/produto, as seções e subseções correspondem à estrutura original de cada documento considerado.\\

{\bf Palavras-chave:} \PalavrasChave.
\end{abstract}
\newpage
\part{Produtos de Fabricio Solagna}
\chapter{Análise de experiências nacionais e internacionais de participação mediada por Internet}
Este documento contempla uma análise de experiências nacionais e internacionais de
participação mediada por Internet considerando os aspectos de inclusão, inovação e
deliberação, com a finalidade de gerar uma matriz de características para integração ao portal
Participa.br. As iniciativas analisadas foram agrupadas em 5 capítulos temáticos, somando 18
iniciativas e mais de 30 ferramentas de participação digital. O capítulo final apresenta
recomendações metodológicas à luz das experiências estudadas.

\section{Introdução}

Página 10

\pp{1} ok.

\pp{2} ok. Unipresente deve ser Omnipresente ou Ubíqua.

\pp{3} ok.

\pp{4} Legal da definição do Macintosh.

\pp{5} Porquê esta visão de participação social se opõe a esta capacidade do indivíduo influenciar os processos? É quanto à visão do Macintosh?

\pp{6} {\bf \color{red} Motivação para o SM essa inspiração através do empírico nas redes sociais}.

\pp{7} Legal a definição de ``gift economics'' como processos nos quais atores cooperam em projetos maiores que as corporações.

\pp{8} Ok.

\noindent Página 11

\pp{9} Dada a função da OGP, é pertinente observar quais os posicionamentos e recursos que exibem quanto à web semântica.

\pp{10} Legal a necessidade da ferramenta de participação de delimitar \emph{input} e \emph{output}.

\pp{11} Muito interessante. Esquematicamente: políticas públicas equilibra: necessidades da comunidade, recursos disponíveis e capacidade de execução. Participação pela internet: há hiato entre capacidade de coletar opinião e sistematização para ação. Congregando colaboração massiva, qualidade e resultado prático, cita: Wikipédia, o Software Livre e o Crowdfunding.

\pp{12} Interessante que nos gov abertos essas iniciativas tiveram pouco impacto.

\noindent Página 12

\pp{13} ok.

\pp{14} Legal da suavização da separação digital e presencial.

\pp{15} ok.

\pp{16} Interessante: inovação, inclusão e protagonismo; deliberação, promoção, construção de políticas públicas.

\pp{17} foco no executivo e na fiscalização de serviço público. Legal do longo histórico com o Legislativo.

\pp{18} ok. Cap 2 -> Democracia digital ligada aa PR nos EUA, Chile, Bolívia e País de Gales.

\noindent Página 13

\pp{19} Cap 3 -> visualização e deliberação de Orçamento Público.
Legal da simulação e visualização pela popularização dos datasets. Seria bom ver quais as iniciativas do OKF caso não esteja neste cap.

\pp{20} Cap4 -> projs/entidades com relevância na democracia digital.

\pp{21} Cap5 -> monitoramento e fiscalização de obras e ações do governo.

\pp{22} Cap6 -> Gabinete Digital do RS. Três maiores consultas nacionais ou internacionais?

\pp{23} Cap7 -> Recomendações de metodologias e ferramentas.

\pp{24} Legal dos três vértices: elaboração, fiscalização e monitoramento.

\noindent Página 14 

\pp{25} Ok.

\section{Experiências de democracia digital ligadas
diretamente a presidência da república}
Página 15
\subsection{Chile: Gobierno Abierto}
Página 16

Figura 1: ok.

\pp{1} ok.

\pp{2} ok.

\pp{3} ok.

\noindent Página 17

Figura 2: ok

\noindent Página 18

\pp{4} ok

\pp{5} ok

\pp{6} ok. Interessante essas consultas que já passaram estarem acessíveis junto a conteúdo jornalístico e multimídia.

\noindent Página 19

Figura 4: ok.

\pp{7} ok.

\noindent Página 20

\subsubsection{Resultados alcançados}

\pp{8} ok.

\subsection{Bolívia: Urna de Cristal}

\indent Figura 5: ok.

\noindent Página 21

\pp{1} ok.

\subsubsection{Metodologia e arquitetura de escolha}

\pp{2} Legal da proposta ou pergunta e então divulgação nas redes sociais.

\pp{3} ok.

\subsubsection{Resultados alcançados}

\pp{4} Interessante o envolvimento das equipes nas redes sociais. {\bf \color{red} Talvez valha eu procurar melhor como fazem isso.}

\noindent Página 22

Figura 6: ok.

\subsection{EUA: Iniciativas de diálogo com o Governo}

\pp{1} ok.

\pp{2} ok.

\pp{3} ok.

\noindent Página 23

\subsubsection{Open for Question}

\pp{4} Legal do envio de questões. Interessante a ambiguidade do texto, que deixou a sugestão de fazer uma relação de pessoas cadastradas para receberem a pergunta.

Figura 7: ok.

Metodologia

\pp{5} bacana a ideia de resposta em video pelo presidente.

\subsubsection{AskObama on Twitter}

\pp{6} super da equipe do Twitter junto.

\noindent Página 24

Figura 8: ok.

\subsubsection{Reddit}

\pp{7} ok.

\noindent Página 25

Figura 9: ok.

\subsubsection{We the People}

\pp{8} Interessante de petição estar integrado ao site da casa branca. {\bf \color{red} Talvez pensar em algo assim para o Participa.br, fornecendo ao usuário meios sofisticados de ativar suas redes sociais.}

Metodologia

\pp{9} ok.

\pp{10} ok.

\pp{11} Legal das respostas quando a petição tem mais de 100 mil assinaturas.

\noindent Página 26

Resultados

\pp{12} ok.

Figura 10: ok.

\subsection{País de Gales: E-petitions}

\pp{13} ok.

\noindent Página 27

Figura 11: ok.

\noindent Página 28

\section{Experiências de visualização e deliberação de Orçamento Público}

\noindent Página 29

\subsection{Orçamento Participativo (OP)}

\pp{1} ok dos três períodos do OP.

\pp{2} Dúbia a construção? Assumindo aqui que a segunda fase é situada entre 1989 e 1992 enquando a primeira fase vai de 1960 até 1998. (?)

\pp{3} OP é criação do PT? (creio que não, mas o texto está me dando a entender isso)

\pp{4} Interessante a mescla de OP com outras iniciativas de democracia digital.

\pp{5} ok.

\pp{6} Bem curiosa essa ausência de OP federal e pouca estadual.

\pp{7} ok.

\noindent Página 30

\subsection{Orçamento Participativo Digital – Belo Horizonte}

\pp{8} ok.

Figura 12: ok.

\pp{9} ok.

\subsubsection{Metodologia do OP Digital}

\pp{10} ok.

\pp{11} ok, bem legal da votação via ligações telefonicas e via app de celular.

\noindent Página 31

\pp{12} ok

Figura 13: ok.

\pp{13} ok.

\subsubsection{Resultados}

\pp{14} ok.

\noindent Página 32

\pp{15} ok.

\subsection{Nova Iorque: Participatory Budgeting in New York City}

\pp{16} ok.

\subsubsection{Metodologia e arquitetura de escolha}

\pp{17} ok.

\noindent Página 33

Figura 14: ok.
Figura 15: ok.

\noindent Página 34

\subsubsection{Orçamento Participativo em outras cidades norte-americanas}

\pp{18} ok.

\pp{19} ok.

\subsection{Liverpool: Budget Simulator}

\pp{20} que tendência?

Figura 16: ok.

\subsection{Stabilize the U.S. Debt}
\noindent Página 35
\pp{21} ok.
\pp{22} ok.
\pp{23} {\bf \color{red} Gamificação}.
\pp{24} ok.

Figura 17: ok.

\noindent Página 36

\subsection{Where Does My Money Go}
\pp{25} ok.
\pp{26} ok.
\pp{27} ok.
\pp{28} ok.
\pp{29} ok.

Figura 18: ok.

\noindent Página 37

\subsubsection{Metodologia utilizada}
\pp{30} ótimo das techs livres.

Figura 19: ok.

\subsection{Aplicações no Brasil}
\subsubsection{Para onde foi meu dinheiro}
ok tudo.
\subsubsection{Orçamento ao seu alcance}
Bem legal sobre o orçamento federal, até porque não há OP federal.

\section{Projetos de democracia digital oriundos da sociedade civil}

\noindent Página 42

\subsection{PortoAlegre.cc}
\pp{1} ok.
\pp{2} ok. Quais conceitos de inteligência social?
\pp{3} ok.
\subsubsection{Metodologia e arquitetura de escolha}
\pp{4} ok.
\pp{5} ok.
\pp{6} Super legal dos pilares: cultura de cidadania, ética do cuidado, corresponsabilização cidadã e engajamento cívico.
\pp{7} 

\noindent Página 43

Figura 23: ok.

\pp{8} ok.
\pp{9} Interessante a confluência de pautas levantadas pela sociedade com agendas públicas.
\pp{10} 

\noindent Página 44

Figura 24: ok.

\subsubsection{Fases de uma causa}
\pp{11} ok.
\pp{12} ok.
\pp{13} ok.
\pp{14} ok.

\subsubsection{Resultados}
\pp{15} ok.
\noindent Página 45
\subsection{MySociet}
\pp{16} ok da UKCOD.
\pp{17} só de transportes?
\pp{18} ok. Legal que trabalham sob encomenda, talvez fazer contato em algum momento ou trazer para uma consultoria presencial.
\pp{19} ok, lembra o lm e orgs de gsoc.
\pp{20} ok.

\subsubsection{Alaveteli}
\pp{21} legal. Em que países são ou foram usados?
\noindent Página 46
\pp{22} ok.
\pp{23} ok.
Figura 25: ok.
\pp{24} ok.
\pp{25} ok.
\noindent Página 47
\subsubsection{FixMyStreet}
\pp{26} ok.
\pp{27} Figura 26: ok.
\pp{28} ok.
\noindent Página 48
\pp{29} ok.
\subsubsection{Eu prometo}
\pp{30} ok. Mandar p alguns amigos.
Figura 27: ok.
\pp{31} ok.
\noindent Página 49
\subsection{SeeClickFix}
\pp{32} ok.
\pp{33} ok.
\pp{34} ok.
Figura 28: ok.
\noindent Página 50
\subsection{All Our Ideas}
\pp{35} ok.
\pp{36} ok.
\pp{37} bem legal. Belo parágrafo sobre All Our Ideas e muito do que a participação social trata no momento.
\subsubsection{Metodologia empregada}
\pp{38} ok.
\noindent Página 51
Figura 29: ok.
\pp{39} ok. Dá vontade de analisar os dados de uma consulta com o Allourideas, também de ver como funciona, e caracterizar estatísticamente. Talvez também propor dinâmicas em rede para melhor coleta de ideias e votos.
\pp{40} ok.
\noindent Página 52
\pp{41} ok. Seria interessante um parâmetro para esse "bastante fiel".
\pp{42} ok.
\pp{43} ok.
\subsection{Cidade Democrática}
\pp{44} ok.
\pp{45} ok. Propor no Cidade Democrática de todo estudante de curso superior público dê 1h de aula gratuita por semana em comunidades menos favorecidas.
\pp{46} ok. Muito legal do sonho p Xingu.
\noindent Página 53
Figura 30: ok.
\noindent Página 54
\section{Ferramentas de monitoramento de obras e políticas públicas}
\noindent Página 55
\pp{1} ok. {\bf \color{red} referência boa para fazer ponte com os avanços recentes e previstos com as redes sociais.}
\pp{2} ok.
\pp{3} ok
\subsection{OnTrack}
\pp{4} ok.
\pp{5} ok.
Figura 31: ok.
Figura 32: ok. Seria bastante interessante saber se há alguma iniciativa que implementa ou visa implementar o OnTrack no Brasil e se podemos usar instâncias já instaladas.
\subsection{De Olho nas Obras}
\pp{6} ok.
\subsection{RapidSMS}
\pp{7} bem legal. {\bf \color{red} Considerando p usar c AA.}
\noindent Página 57
\pp{8} Muito bom. Seria bastante interessante saber se há alguma iniciativa que implementa ou visa implementar o OnTrack no Brasil e se podemos usar instâncias já instaladas.
Figura 33: ok.
\subsection{Colab}
\pp{9} bem interessante. {\bf \color{red} ver sobre as ferramentas e entrar em contato com equipe}.
\noindent Página 58
\section{Estudo de caso: Gabinete Digital do Rio Grande do Sul}
\noindent Página 59
\pp{1} ok.
\pp{2} ok.
\pp{3} ok.
\pp{4} ok.
\pp{5} ok.
\pp{6} ok. {\bf \color{red} procurar e navegar minimamente pelo repo}.
\pp{7} ok.
\noindent Página 60
\subsection{Ferramentas de participação do Gabinete Digital RS}
\subsection{Governador Pergunta}
\pp{8} ok.
\pp{9} muito legal. {\bf \color{red} Temos acesso aos dados desta(s) consulta(s)?}
\pp{10} ok. {\bf \color{red} Bom que esta grande parcela da população reforça a pertinência do uso destes dados para benchmarks e pesquisas científicas.}
\pp{11} Interessante.
\subsubsection{Primeira edição: saúde pública}
\pp{12} muito legal.
\noindent Página 61
Figura 34: ok
\subsubsection{Segunda edição: Segurança no Trânsito}
\pp{13} ok.
\noindent Página 62
Figura 35: ok
\subsubsection{Terceira Edição: Reforma política}
\pp{14} ok.
\pp{15} ok.
\pp{16} ok.
\pp{17} ok. {\bf \color{red} Há opção no PairWise para equivalência entre as ideias? (assim o sistema pode estabelecer propostas equivalentes).}
\pp{18} ok.
\noindent Página 63
\pp{19} ok. Seria interessante fazer plots:
\begin{itemize}
    \item tempo de consulta \emph{versus} número de votos. 
    \item tempo de consulta \emph{versus} número de ideias.
    \item número de votos \emph{versus} número de ideias.
\end{itemize}
\subsubsection{Passos e condições para as consultas}
\pp{20} ok.
\pp{21} ok, muito legal.
\pp{22} ok, muito legal. {\bf \color{red} há ou foi discutida uma página similar para acompanhamento das propostas do arenaNETmundial?}
\noindent Página 64
Figura 36: ok.
\pp{23} ok.
\pp{24} ok.
\noindent Página 65
\pp{25} ok.
\subsubsection{Metodologia e arquitetura de escolha empregadas no Governo Pergunta}
\pp{26} ok.
\pp{27} ok.
\pp{28} ok, muito legal desse código. {\bf \color{red} procurar o repo e navegar minimamente.}
\pp{29} ok. Essa implementação evitou a necessidade de instalar várias instâncias do AllOurIdeas (pairwise?), é isso?
\noindent Página 66
Figura 37: ok.
\subsubsection{Propostas sem mediação}
\pp{30} ok.
\pp{31} ok.
\pp{32} ok.
\noindent Página 67
\pp{33} ok. {\bf \color{red} Mesmo positivo neutralizar grupos de interesse?}.
\pp{34} ok.
\pp{35} Interessante saber que as propostas polêmicas não tem força neste contexto. {\bf \color{red} elaborar formas de equilibrar esse processo?}
Tabela 1: ok.
\pp{36} ok.
\noindent Página 68
\subsubsection{Mobilização e Vans de Participação}
\pp{37} ok.
Figura 38: ok.
\pp{38} ok.
\pp{39} ok.
\noindent Página 69
\pp{40} ok.
\pp{41} ok.
\pp{42} ok.
\subsection{Governador Responde}
\pp{43} ok.
\pp{44} ok.
\pp{45} ok.
\noindent Página 70
\pp{46} ok
Tabela 2: ok.
\pp{47} ok.
\pp{48} ok. Uma pena estar desativada. {\bf \color{red} Seria legal ter acesso a estes dados}
\noindent Página 71
Tabela 3: ok.
Figura 39: ok.
\noindent Página 72
\subsection{Governo Escuta}
\pp{49} ok.
\pp{50} ok.
Tabela 4: ok.
\pp{51} ok.
\pp{52} ok.
\pp{53} ok.
\noindent Página 73
Figura 40: ok.
\subsubsection{Diálogo com 500 mil pessoas}
\pp{54} ok.
\pp{55} {\bf \color{red} Interessante que o número de pessoas nas ruas foi aprox. igual ao de pessoas assistindo do Governador escuta. Interessante que das 500 mil visitações, aprox. 4\% assistiram ao video, que é a parcela esperada de hubs.}
\pp{56} ok.
\subsection{Agenda Colaborativa}
\pp{57} ok.
\noindent Página 74
\pp{58} ok.
Figura 41: ok.
\pp{59} ok
\noindent Página 75
\pp{60} ok.
Tabela 5: ok.
\pp{61} ok.
\pp{62} ok, bem legal estas oficinas.
\pp{63} ok.
\subsection{Monitoramento de obras públicas}
\pp{64} ok.
\pp{65} ok.
Figura: 42: ok.
\noindent Página 77
\subsubsection{Metodologia}
\pp{66} ok.
\pp{67} ok.
\pp{68} ok.
\pp{69} ok.
\subsection{Prêmios alcançados}
\pp{70} ok.
\noindent Página 78
\subsection{Sistema Estadual de Paritcipação Cidadã}
\pp{71} ok.
\pp{72} ok. Legal do site da participação.

\noindent Página 79
\section{Proposta de integração e metodologias para o Participa.br}
\noindent Página 80
\pp{1} ok.
\pp{2} ok.
\pp{3} ok.
\pp{4} ok.
\subsection{Instâncias, mecanismos e ferramentas de participação}
\pp{5} ok.
\pp{6} ok.
\pp{7} ok.
\noindent Página 81
Figura 43: ok. Bem legal.
\pp{8} ok.
\pp{9} ok.
\pp{10} ok. {\bf \color{red} dicernimento útil: reflexivo, não imperativo.}
\pp{11} ok.
\noindent Página 82
Figura 44: ok.
\subsection{Fases das políticas públicas e a participação social}
\pp{12} ok.
\pp{13} ok.
\pp{14} ok. {\bf \color{red} Interessante essa relação entre bottom-up, democracia participativa e legitimidade de política pública.}
\pp{15} {\bf \color{red} Instrumental esta divisão do sebrae.}
\noindent Página 83
\pp{16} ok.
\subsection{Metodologias para mecanismos de participação}
\pp{17} ok.
\pp{18} ok.
\noindent Figura 45: ok.
\pp{19} ok.
\pp{20} ok. {\bf \color{red} Precendente para a foco na categorização dos agentes participativos: citada como um desafio pelo consultor.}
\noindent Página 84
\pp{21} ok.
\subsection{Arquitetura de escolha}
\pp{22} ok.
\pp{23} ok.
\pp{24} ok.
\pp{25} ok. {\bf \color{red} Talvez passar um período lendo sobre nudges \url{http://www.inudgeyou.com}.}
\noindent Página 85
\pp{26} ok. Talvez visitar o livro.
\pp{27} ok. {\bf \color{red} Como pode refletir positivamente: através de uma cartilha dos métodos mais comuns, para instrumentalizar o participante a reconhecer caso haja manipulação indevida.}
\pp{28} ok.
\pp{29} ok.
\pp{30} ok.
\noindent Página 86
\pp{31} ok.
\pp{34} ok.
\subsection{Mecanismos e metodologias recomendados}
\subsubsection{Constituição da agenda}
\pp{35} ok.
\pp{36} ok.
\pp{37} ok.
\pp{38} ok.
\noindent Página 87
\pp{39} ok.
\pp{40} ok.
\pp{41} ok.
\pp{42} ok.
Figura 46: ok.
\noindent Página 88
\subsubsection{Consultas e audiências públicas}
\pp{43} ok.
\pp{44} ok.
\pp{45} ok.
\pp{46} ok. {\bf \color{red} Apontamento para enquetes. Útil caso sejam propostas enquetes que atentem para a diversidade de papéis nas redes sociais.}
\pp{47} ok.
\pp{48} ok.
\pp{49} ok.
\noindent Página 89
Figura 47: ok.
\pp{50} ok.
\subsubsection{Conferências e gestão para conselhos}
\pp{51} ok.
\pp{52} ok.
\pp{53} ok. Tem também as conferências livres de educação, embora n incorporadas ao CONAE como esperado.
\pp{54} ok.
\noindent Página 90
\pp{55} ok.
Figura 48: ok.
\subsubsection{Decisões orçamentárias}
\pp{56} ok.
\pp{57} ok.
\pp{58} ok. {\bf \color{red} ver a ferramenta Delib}.
\noindent Página 91
\subsubsection{Monitoramento de políticas públicas}
\pp{59} ok.
\pp{60} ok.
\pp{61} ok.
\pp{62} ok.
\pp{63} ok.
\pp{64} ok. {\bf \color{red} Na definição usada pelo consultor, ``nós'' são os muito conectados, e os ``hubs'' são os atravessadores. Esta nomenclatura conflita com a da ciência das redes na física, na matemática aplicada e na computação, em que ``hubs'' são os muito conectados e ``nós'' é sinônimo de ``vértices'' (usualmente agentes em redes sociais ou elementos sendo relacionados nestas e outras redes). É útil procurar as fontes do consultor para complementar o ponto de vista usado pela consultoria de dados linkados, redes complexas e processamento de linguagem natural.}
\pp{65} ok.
\noindent Página 92
\pp{66} ok. {\bf Precedente para o uso de critérios de redes para seleção de agentes (p.ex. delegados): o consultor aponta a pertinência destas métricas para complementar os votos.}
\section{Conclusões preliminares}
\pp{1} ok.
\pp{2} ok.
\pp{3} ok.

\chapter{Metodologias de participação, ciclo de políticas públicas e  produção de metadados: integração às trilhas de participação do portal Participa.br}
Este documento contempla uma série de propostas de metodologias de participação mediadas por internet, considerando o ciclo de políticas públicas e sua consequente produção de metadados. O intuito é auxiliar na identificação dos processos decisórios e de participação a partir da plataforma do Participa.br. A proposta também aponta para formas de integração dos processos participativos às chamadas trilhas de participação do portal.
\noindent Página 7
\section{Introdução}
\pp{1} ok.
\pp{2} ok.
\pp{3} ok.
\pp{4} ok.
\pp{5} ok.
\pp{6} ok.
\noindent Página 8
\pp{7} ok. Legal que houve integração com os servidores.
\pp{8} ok.
\noindent Página 9
\section{Ciclo de políticas públicas}
\pp{9} {\bf \color{red} bem legal o triangulo: ciclo de políticas públicas, sistema de participação e trilha participativa.}
\pp{10} ok.
\pp{11} ok.
\pp{12} ok.
\pp{13} que accountability?
\pp{14} ok.
\pp{15} {\bf \color{red} legal as tendências a profissionalização e rarefação dos agentes participativos.}
\noindent Página 10
\pp{16} ok.
Figura 1: Conselhos são prioritariamente para monitoramento e fiscalização? E ouvidoria?
\pp{17} ok.
\pp{18} ok. Interessante dos conselhos e ouvidorias.
\noindent Página 11
\pp{19} ok. {\bf \color{red} ``fugir do setorialismo sem tornar as pautas individualizadas''}
\pp{20} ok.
\pp{21} ok.
\pp{22} ok.
\pp{23} legal dos focos de atenção pelos filtros participativos.
\pp{24} ``consolidar agenda para movimentos sociais''.
\noindent Página 12
\pp{25} dualidade do participa.br: instância de comunidade formal ou por interesse, instância para aflorar demandas individuais pelos processos.
\pp{26} ok. {\bf \color{red} indicador de inclusão na participação social.}
\pp{27} ok.
\pp{28} ok. {\bf \color{red} Declaração de registro histórico com a consulta do arenaNETmundial: 280 mil votos, 200 participações; uma das maiores consultas públicas feitas no Brasil, certamente a maior consulta digital. (internacional?)}
\pp{29} {\bf \color{red} participação digital: novos métodos e melhoras para os existentes}.
\subsection{Trilhas de participação}
\pp{30} ok.
\noindent Página 13
\pp{31} ok.
\pp{32} ok.
\pp{33} ok.
\noindent Página 14
Figura 2: ok.
\noindent Página 15
\pp{34} ok.
\noindent Página 16
Figura 3: ok.
\pp{35} ok.
\pp{36} ok. {\bf \color{red} Pensando em trilhas p o servidor}
\pp{37} ok.
\pp{38} ok.
\pp{39} ok. {\bf \color{red} Apontanda demanda por sistematização dos conteúdos de participação}
\noindent Página 17
\pp{40} ok.
\pp{41} {\bf \color{red} motivação para a sistematização dos conteúdos: evitar sobreposição de pautas e demandas}.
\pp{42} ``balcanização da participação''.
\pp{43} ok.
\pp{44} {\bf \color{red} triângulo: Noosfero incorpora ferramentas participativas; informacionalmente, integra-se com os ciclos participativos e das políticas públicas; metodologicamente, deve ser um canal de uso das ferramentas para as comunidades.}
\noindent Página 18
\subsection{Recomendações para integração das trilhas aos ciclos}
\pp{45} ok.
Figura 4: ok.
\pp{46} ok.
\noindent Página 19
\section{Integração do Participa.br com as mesas de diálogo e monitoramento}
\pp{47} ``interface para mesas de diálogo e monitoramento.''
\pp{48} ok.
\pp{49} ok.
\pp{50} ok. Bem legal das mesas de diálogo (negociação) e de monitoramento.
\pp{51} ok. {\bf \color{red} Demanda de um dashboards de indicadores, para mesas de monitoramento.}
\section{Elementos motivadores}
\pp{52} ok. {\bf \color{red} governança compartilhada através da transparência, prestação de contas e siálogo social}
\pp{53} {\bf \color{red} objs da SNAS: canais participativos não hierarquizados, flexíveis, abertos e que dialógue com as redes.}
\pp{54} O setor está em fase avançada de implementação de um sistema de indicadores, painel de monitoramento e controle de demandas que pode ser integrado ao Partipa.br. {\bf \color{red} onde está este painel? há especificações deste sistema de indicadores?}
\pp{55} {\bf \color{red} ``...auxiliado por processos de dados abertos, provido por uma interface robusta no que tange a aplicação de recursos federais nos diversos níveis, acessível ao cidadão.''}
\pp{56} ok.
\subsection{Recomendações metodológicas}
\pp{57} ``trilha de participação e monitoramento''.
\subsubsection{Estabelecer critérios e requisitos mínimos para a formação de comunidades}
\pp{58} {\bf \color{red} estabelecer claramente que o servidor possui uma trilha de monitoramento que tem como objetivo monitorar o Participa.br. A ``abertura da abertura''.}
\pp{59} ok.
\pp{60} ok.
\subsubsection{Estabelecimento de agenda de participação}
\pp{61} {\bf \color{red} pensar em séries temporais para agendas, além disso, somente os vinculados à SNAS?}
\pp{62} ok.
\noindent Página 22
\pp{63} ok.
\subsubsection{Estabelecer processos dialógicos com prazo definido de retornos}
\pp{64} {\bf \color{red} apontando necessidade de priorizar devolutivas do governo na agenda para melhorar engajamento}.
\pp{65} ok.
\pp{66} ok. {\bf \color{red} procurar os produtos desta consultoria feita para o Departamento de Diálogos Sociais.}
\noindent Página 23
Figura 24: ok. Bastante interessante o ciclo.
\noindent Página 24
\pp{67} {\bf \color{red} proposto sistema de monitoramento para uso interno da SNAS, enquanto o Participa.br p contato com a sociedade.} (isso?)
\pp{68} ok.
\pp{69} {\bf \color{red} apontada pertinência de auxílio na devolutiva do Estado através dos rastros digitais e análise pública deles}.
\subsubsection{Utilização de dados abertos como forma de alimentar indicadores}
\pp{70} {\bf \color{red} Argumento para dados linkados.}
\noindent Página 25
\section{Sugestão para criação de metadados}
\pp{1} ok.
\pp{2} ok.
\pp{3} {\bf \color{red} como assim } ``complementa com mais três além de adicionar outros campos. ''?
\pp{4} ok. (com tabela sem numeração)
\noindent Página 26
\pp{5} legal essa de dc+mpeg 7.
\pp{6} ok.
\subsection{O que classificar?}
\pp{7} ok.
\pp{8} ok. {\bf \color{red} Proposta de metadados para os resultados, trilhas e comunidades}
\subsection{O que é obrigatório e opcional}
\pp{9} {\bf \color{red} como assim?} ``Os campos não devem ter variação quanto à categoria e a tipo dos dados a serem preenchidos''.
\noindent Página 27
\subsection{Vocabulário}
\pp{10} ok.
\subsection{Construindo uma proposta de metadados}
\pp{11} importante {\bf \color{red} apontada pertinência de uma atividade com desenvolvedores, consultores coordenadores e outros envolvidos no projeto, para atividade volta p vocabulário}
\pp{12} legal dos campos para integração com políticas públicas.
Tabela: {\bf \color{red} REFERENCIA: proposta de campos para trilhas participativas.} Integrar em versão próxima da triplificação dos dados do participa.
\noindent Página 28
\pp{13} {\bf pensar na compatibilização da triplificação feita, proposta nesta tabela, e microdados feitos pela Dani.}
\noindent Página 29
\subsection{Integração de metadados, trilhas, ciclos e processos decisórios}
\pp{14} ok.
\pp{15} ok. {\bf \color{red} ocorre-me de fazer um servidor p ação: trabalhar na opa e extração de info dos dados do Participa}
\pp{16} ok.
\pp{17} ok. {\bf \color{red} que grupo de foco?}
\noindent Página 30
\section{Conclusões preliminares}
\pp{1} ok.
\pp{2} ok.
\pp{3} ok. {\bf \color{red} aqui tá ``uma das maiores consultas'' para a do arenaNETmundial}.
\pp{4} ok.
\pp{5} ok. {\bf \color{red} apologia ao sistma de monitoramento encaminhado à DITEC}.


\part{Produtos de Grazielle Machado Fernandes}
\chapter{Plano de trabalho contendo estratégias de relacionamento e mobilização e mapeamento de intelocutores do ``Projeto PARTICIPA''}
\noindent Página 1
\section{Participa.br - Portal Federal de Participação Social}
\pp{1} ok.
\pp{2} ok.
\pp{3} ok.
\pp{4} ok. {\bf \color{red} Necessidade do sistema monitoramento}
\noindent Página 2
Figura: ok. {\bf \color{red} ciclo fundamental encontrado.}
\section{Monitoramento e mapeamento da redes}
\pp{5} ok.
\pp{6} ok.
\noindent Página 3
\section{Identificação de Influenciadores X Gestores Públicos}
\pp{7} {\bf \color{red} previsto SM e contato com agentes detectados e considerados estratégicos}
\pp{8} essa abertura de dados e análise vinga como ideia pois está ancorada na realidade, portanto é útil e instrumental. {\bf \color{red} este momento aponta para o monitoramento interno.}
Figura: ok. {\bf \color{red} mais um ciclo fundamental.}
\section{Mapeamento e abordagem de interlocutores}
\pp{9} ok.
\noindent Página 4
\pp{10} {\bf \color{red} Trecho descreve trilhas criadas via colaboração de agentes civis (escuta/monitoramento nas redes sociais) e gov (monitoramento interno)}.
\pp{11} ok e previsto como responsabilidade da equipe do Participa.br.
\pp{12} ok. {\bf \color{red} estas regras são campos já presentes nas trilhas? Colocar na OPA}
\pp{13} ok.
\section{Ações previstas para realização até dezembro de 2014}
\pp{14} ok.
\pp{15} ok.
\pp{16} ok.
\pp{17} ok. {\bf \color{red} para a difusão do Participa.br e termo \#ParticipaBR, escolher agentes para fazerem as 3 rodadas de difusão nas redes deles, ou comunicação única através de atravessadores e mais próximos.}
\noindent Página 5
\section{Ações a serem executadas para atigir os objetivos propostos}
\pp{18}




\part{Produtos de Joênio Costa}

\part{Notas aleatórias}

\chapter{Simulação de processos}

Difusão, contágio, acessibilidade e fofoca nas redes reais do Participa.br e de outras instâncias.

\chapter{Criação de benchmarks participativos}
Com medidas de experiências participativas (quantidade de participantes, médias, interações, etc)

\chapter{Triplificação dos dados participativos}
Recomendar p equipe de dev do noosfero que gerem URLs para acessar comentários específicos e outros recursos em conformidade com as URIs usadas na triplificação dos dados do Participa.br.

Triplificar os dados do Cidade Democrática e do Participatório.

Triplificar os dados da consulta do arenaNETmundial.

Triplificar os dados das experiências do RS.

Triplificar Tweets participativos meus, do Malini do Pimentel e de outros.


\newpage
\bibliography{bibliografia}
\newpage
%\listoffigures
\section*{Abreviações e jargão}
\begin{itemize}[label={}]
    \item {\bf RC:              } Redes Complexas
    \item {\bf PLN:             } Processamento de Linguagem Natural
    \item {\bf OPS:             } Ontologia de participação Social
    \item {\bf OPA:             } Ontologia do Participa.br
    \item {\bf MMISSA:          } Monitoramento Massivo e Interativo da Sociedade pela Sociedade para Aproveitamento
    \item {\bf AARS:            } A Análise de Redes Sociais
    \item {\bf MyNSA:           } Monitoring yields Natural Streaming and Analysis
    \item {\bf PNPS:            } Plano Nacional de Participação Social
    \item {\bf RDF:             } Resource Description Framework
    \item {\bf HTTP:            } Hypertext Transfer Protocol
    \item {\bf SPARQL:          } Simple Protocol and RDF Query Language
    \item {\bf endpoint SPARQL: } ponto de acesso, geralmente HTTP, a dados em RDF via buscas em SPARQL.
    \item {\bf Participa.br:    } Portal federal de participação social.
    \item {\bf IPython Notebook:} instância online para rodar scritps Python
    \item {\bf Mateor:          } arcabouço para páginas reativas e com funcionamento distribuído.
    \item {\bf D3js:            } biblioteca de visualização de dados.
\end{itemize}

\newpage
\printindex
\newpage
%

\appendix
\end{document}
